\section{Introduction}
\label{sec:intro}

Periodic composites modify wave propagation through their band structure. Early calculations and experiments established that two-dimensional elastic composites can exhibit complete acoustic gaps \cite{SigalasEconomou1992,kushwaha1993,vasseur2001}. Subsequent lattice calculations and topology optimization addressed the placement of pass bands and stop bands \cite{MartinssonMovchan2003,SigmundJensen2003,Dalklint2022}. More recent work has examined frequency-targeted gaps, multimaterial designs and their dependence on structural size \cite{He2026,Nakayama2025,Quinteros2026}. Broader accounts of phononic crystals and their band-gap mechanisms are available \cite{Vasileiadis2021,Oudich2023,MaWangYuan2026}. A frequency interval free of modes along one wave-vector path may still contain modes elsewhere in the Brillouin zone. Recent studies have considered active Bloch-wave control, resonant metadamping and higher-gradient beam models with micro-inertia \cite{IJESBadinoBacigalupo2026,IJESZhaoZhangWuHu2026,IJESAyadKarathanasopoulosGanghofferLakiss2020}. Further work examines material length scales in peridynamics and the relation between different strain-gradient theories \cite{IJESChanChen2023,IJESFuZhouQi2020}.

When a wavelength approaches the material's microstructural scale, a local Cauchy law may not capture size-dependent stiffness or dispersion. Higher-order descriptions account for gradients of deformation and the stresses conjugate to them. The relevant foundations include couple-stress and microstructure theories, second-gradient elasticity and higher-order boundary work \cite{toupin1962,MindlinTiersten1962,mindlin1964,Mindlin1965,mindlineshel1968,Germain1973}. Later studies examined strain-gradient descriptions of nonuniform deformation and small-scale beam experiments \cite{Aifantis1999,YangChongLamTong2002,LamYangChongWangTong2003}. Reviews and lattice-based derivations describe other formulations of gradient elasticity \cite{askesaifantis2011,polyzosfotiadis2012}. Peridynamic wave-dispersion studies have also compared nonlocal kernels with other nonlocal elasticity models \cite{IJESEremeyevNaumenko2025}. We use the local Mindlin Form-II theory; nonlocal integral and micromorphic theories are related approaches, but they are distinct formulations.

A dynamic gradient model needs a kinetic-energy law as well as gradient stiffness. With classical inertia alone, the shear branch examined here has phase speed proportional to wavenumber at high wavenumbers. Independent gradient micro-inertia gives this branch a finite high-wavenumber limit. That result follows from the constitutive and kinetic-energy laws used here; other higher-order continua need not have the same limit. One-dimensional studies have examined dispersion in gradient-elastic beams and periodic media, including dipolar-gradient formulations \cite{papargyribeskou2009,liweizhou2016,HosseiniSladekSladekZhang2024}. A two-dimensional phononic crystal has also been analyzed with a relaxed micromorphic model that includes free and gradient micro-inertia \cite{MadeoColletMiniaciBillonOuisseNeff2018}. Nonlocal strain-gradient formulations have also been used to examine how material length scales affect dispersion and group velocity \cite{IJESLazarAgiasofitou2026}. These formulations differ in their treatment of inertia and higher-order stress. The calculations below use local Mindlin Form-II elasticity.

Several constitutive descriptions have been used to study size effects in phononic crystals. Transfer-matrix analyses show that imperfect interfaces can alter one-dimensional band gaps \cite{zhengwei2009}. Other one-dimensional work considers dipolar-gradient elasticity, gradient thermoelasticity, and flexoelectric strain-gradient coupling \cite{liweizhou2016,li2023anchorA,li2024anchorB,HosseiniSladekSladekZhang2024}. In two dimensions, researchers have applied nonlocal strain-gradient theory to nanoscale phononic crystals and studied mixed and defect modes in a phononic-crystal slab \cite{JinHuHu2022,JinHuHu2023}. Other model-specific comparisons include thermoelastic cylindrical crystals and functionally graded lattices \cite{hosseinizhang2021,mishra2026anchorC}. Anisotropy and orientation have also been studied in conventional phononic crystals: work on orthotropic fillers and anisotropic plates reports shifts in band locations and widths as material axes or inclusions rotate relative to the lattice \cite{zhanwei2010,YaoWuHouXin2011}. Recent studies examine direction-dependent gaps in anisotropic nanocrystal assemblies and the effects of anisotropic constituents in optimized phononic crystals \cite{Kim2023,Liu2024anisotropic}. For flexural plate waves, anisotropic equal-frequency contours can align group velocity with a lens interface over a broad frequency range \cite{Danawe2023}. These studies concern the orientation of inclusions or elastic anisotropy. Here the orientation is assigned to the characteristic-length tensor itself, with gradient micro-inertia included in the kinetic energy.

The strain-gradient weak form contains second derivatives of displacement. A displacement-based approximation therefore needs an $H^2$-conforming trial space; mixed formulations provide another option. The bicubic Bogner--Fox--Schmit (BFS) rectangle provides $C^1$ continuity across elements on structured quadrilateral meshes \cite{bfs1965}. Earlier studies developed displacement-based, mixed, and implicit-gradient finite elements \cite{ShuKingFleck1999,AmanatidouAravas2002,ZervosPapanicolopulosVardoulakis2009,Beheshti2017}. Later work treats second-gradient boundary-value problems, higher-order Hermite mappings on regular and distorted domains, mixed elements with patch recovery or shared degrees of freedom, and mixed crack formulations \cite{AndreausDellIsolaGiorgioPlacidiLekszyckiRizzi2016,BacciocchiFantuzzi2023,BacciocchiFantuzzi2025,ChoiLeeSim2023,Papanicolopulos2024,ChirkovNazarenkoAltenbach2024}. For a conforming BFS Bloch element, the same phase transformation must be applied to displacement values and all derivative degrees of freedom, including the mixed derivative.

We consider a two-dimensional Bloch problem with an oriented ellipsoidal length tensor, independent gradient micro-inertia and conforming $C^1$ interpolation. The numerical study has four components:
\begin{itemize}
  \item We formulate a Bloch--Floquet framework for two-dimensional anisotropic strain-gradient elasticity. The rotated ellipsoidal averaging domain defines the length tensor, with aspect ratio $\mathrm{AR}=l_1/l_2$ and orientation $\theta$ as separate design variables.
  \item Gradient micro-inertia is examined analytically and numerically. For the branch studied, $\ell_{\mathrm{i}}=0$ gives an unbounded high-wavenumber phase speed, whereas $\ell_{\mathrm{i}}>0$ yields the finite limit $v_{T,\infty}=\sqrt{\mu/\rho}\,l_{\mathrm{eff}}/(\sqrt{10}\,\ell_{\mathrm{i}})$ derived in Appendix~\ref{app:asymptotics}.
  \item A conforming 32-degree-of-freedom BFS element is developed, with the same complex Bloch phase applied to displacement, first derivatives, and the mixed derivative.
  \item Across the $(\theta,\mathrm{AR})$ design space, we evaluate directional dispersion, gap measures, and group velocity. The energy-flux relation follows from local balance; the analysis also separates directional spectral separations from complete gaps and measures steering amplitude and lobe orientation.
\end{itemize}

{\sloppy
For the homogeneous case, the Bloch cell is a computational construction; it contains no material contrast. A separation confined to part of the reduced-zone path is thus not a Bragg attenuation gap. This distinction has been discussed for anisotropic media \cite{WangMaznevLaude2016}. Earlier analyses of periodic lattices have also treated directional beam steering using homogenized continua, including a hexagonal-lattice model \cite{RosiAuffray2019}. Our homogeneous calculations report reduced-zone separations and group-velocity steering. The periodic-composite calculation is treated separately because its band gap arises in a materially contrasted cell.
\par}

